\chapter{Turunan: Satu dan Banyak Variabel}

Sebuah pencatat posisi pada sepeda menulis satu titik setiap detik. Untuk memperoleh
kecepatan, kita membagi selisih jarak dengan selisih waktu. Dua bacaan berjarak satu detik
dan terpisah empat meter memberi empat meter per detik, meskipun itu hanya kecepatan
rata-rata. Bila jaraknya seperseratus detik, selisih jaraknya beberapa sentimeter sedangkan
derau alat beberapa meter; hasil baginya melompat ke ratusan meter per detik, bahkan bertanda
negatif. Pembagian yang sama, dua jawaban yang berbeda.

Turunan menyelesaikan ketegangan itu dengan memindahkan pertanyaan ke limit, tempat derau
tidak masuk hitungan. Tetapi limit tidak dapat diukur: yang terukur hanya hasil bagi pada
jarak berhingga, dan ia membawa dua galat sekaligus, pemotongan karena jaraknya tidak nol
dan pembulatan karena mesin hanya menyimpan sekian digit. Turunan adalah objek matematis
yang tepat; hasil bagi adalah aproksimasinya, dengan galat yang dapat dihitung dan diatur.

Kesulitan kedua muncul ketika masukan lebih dari satu. Sebuah model penilaian kredit
menerima dua belas masukan, dan menggeser satu masukan memberi satu angka laju perubahan.
Menggeser dua masukan sekaligus tidak sama dengan menjumlahkan kedua angka itu, kecuali
fungsinya berperilaku baik; nanti kita lihat contoh yang gagal, dan fungsi itu tetap
kontinu. Jawaban yang benar bukan daftar dua belas angka, melainkan satu peta linear pada vektor
pergeseran: itulah turunan sesungguhnya, dan matriksnya Jacobian. Kekontinuan Bab 11 tidak
mengatakan seberapa cepat sesuatu berubah.

\begin{untukanak}
Pernahkah kamu berdiri di atas atap yang melengkung? Kalau kamu melangkah ke depan, kamu
mendaki. Kalau ke samping, kamu bisa tetap pada ketinggian yang sama. Padahal kamu berdiri di
titik yang sama, di atap yang sama. Jadi pertanyaan ``seberapa curam tempat ini'' belum
lengkap sebelum kita menyebut arah langkah.

Ada masalah lain. Kalau kamu hanya melihat dua titik di atap, lalu membagi selisih ketinggian
dengan selisih jarak, kamu mendapat kemiringan rata-rata, bukan kemiringan di satu titik.
Untuk mendapat kemiringan tepat di bawah kakimu, kedua titik harus didekatkan terus-menerus.
Semakin dekat semakin tepat, tetapi kalau terlalu dekat angka pengukuranmu kacau karena alat
ukur tidak sempurna.

Di bab ini kita akan tahu mengapa kemiringan di satu titik butuh arah, mengapa arah paling
curam selalu ada, dan mengapa alat hitung yang terlalu rapat jaraknya justru salah.
\end{untukanak}

\section{Turunan sebagai limit dan dua artinya}
\tingkat{dasar}

Laju perubahan sesaat adalah bilangan yang didekati hasil bagi dua selisih ketika selisih
masukannya mengecil. Definisi berikut hanya menuliskan kalimat itu dengan hati-hati, karena
hasil baginya sendiri tidak terdefinisi ketika penyebutnya nol.

\kamus{turunan}{derivative}{Kemiringan garis yang paling rapat menempel pada grafik di satu
titik, yang sekaligus menjadi laju perubahan sesaat di titik itu.}

\begin{definisi}[Turunan di satu titik]
Misalkan $f$ terdefinisi pada interval terbuka yang memuat $c$. Turunan $f$ di $c$ adalah
\[
  f'(c)=\lim_{h\to0}\frac{f(c+h)-f(c)}{h},
\]
asalkan limit itu ada dan berhingga. Dalam hal ini $f$ disebut
\emph{terdiferensialkan} di $c$. Notasi lain yang dipakai di buku ini: $Df(c)$ dan
$\dfrac{df}{dx}(c)$.\index{turunan}
\end{definisi}

Hasil bagi itu terdefinisi untuk setiap $h\neq0$, jadi tidak ada pembagian dengan nol; yang
mungkin gagal adalah limitnya. Kegagalan itu nyata: untuk $f(x)=|x|$ di $c=0$ hasil baginya
$1$ untuk $h>0$ dan $-1$ untuk $h<0$, sehingga limit duanya tidak ada.

Secara fisis, turunan adalah laju perubahan sesaat: $s'(t)$ kecepatan bila $s(t)$ posisi,
$Q'(t)$ arus bila $Q(t)$ muatan yang mengalir, dan $C'(q)$ biaya marginal bila $C(q)$ biaya
produksi $q$ satuan.

Secara geometris, hasil bagi adalah kemiringan talibusur dari $(c,f(c))$ ke
$(c+h,f(c+h))$. Ketika $h\to0$, talibusur itu berputar mengelilingi titik pertama menuju satu
posisi; garis pada posisi limit disebut garis singgung, dan kemiringannya $f'(c)$. Persamaan
garis singgung di $c$ adalah
\[
  y=f(c)+f'(c)(x-c).
\]
Garis singgung bukan ``garis yang menyentuh grafik tepat satu kali'', sebab untuk
$f(x)=x^{3}$ di $c=0$ garis singgung $y=0$ memotong grafik di titik asal. Yang benar:
selisihnya terhadap grafik mengecil lebih cepat daripada jaraknya.

Putaran itu dapat diikuti angkanya. Untuk $f(x)=x^{2}$ di $c=2$, talibusur dengan $h=1$
berkemiringan $5$, dengan $h=0{,}5$ kemiringan $4{,}5$, dan garis singgungnya $4$.
Gambar~\ref{fig:sekan-singgung} menggambar ketiganya: dari $5$ ke $4{,}5$, menetap di $4$.

\begin{figure}[htbp]
\centering
\begin{tikzpicture}
\begin{axis}[
  width=\linewidth, height=6.4cm,
  axis lines=left,
  xlabel={$x$}, ylabel={$y$},
  xmin=1.55, xmax=3.25, ymin=2.2, ymax=10.2,
  xtick={2,3}, ytick={4,9},
  tick label style={font=\footnotesize},
  label style={font=\footnotesize},
  legend style={font=\scriptsize, draw=none, fill=none,
                at={(0.03,0.97)}, anchor=north west},
  clip=false,
]
\addplot[ink, thick, domain=1.6:3.15, samples=90] {x^2};
\addlegendentry{$f(x)=x^{2}$}
\addplot[accent, thick, domain=1.95:3.18] {5*(x-2)+4};
\addlegendentry{talibusur $h=1$, kemiringan $5$}
\addplot[layerteal, thick, dashed, domain=1.95:3.0] {4.5*(x-2)+4};
\addlegendentry{talibusur $h=0{,}5$, kemiringan $4{,}5$}
\addplot[layerindigo, thick, domain=1.68:3.05] {4*(x-2)+4};
\addlegendentry{garis singgung, kemiringan $4$}
\addplot[only marks, mark=*, mark size=1.5pt, ink]
  coordinates {(2,4) (2.5,6.25) (3,9)};
\node[font=\scriptsize, anchor=north west, text=ink] at (axis cs:2.04,3.75) {$(2,4)$};
\end{axis}
\end{tikzpicture}
\caption{Talibusur yang berputar menjadi garis singgung pada $f(x)=x^{2}$ di $(2,4)$.
Kemiringannya turun dari $5$ (untuk $h=1$) ke $4{,}5$ (untuk $h=0{,}5$) dan menetap di $4$
ketika $h\to0$.}
\label{fig:sekan-singgung}
\end{figure}

\begin{bacakeras}
Dua suara bergantian; jawabannya makin pendek setiap kali.

--- Berapa cepat dia bergerak? --- Dari titik mana ke titik mana?

--- Dari detik pertama ke detik kedua. --- Tiga meter.

--- Dari detik pertama ke detik satu setengah. --- Dua setengah meter.

--- Dari detik pertama ke detik satu sepersepuluh. --- Dua satu per lima.

--- Dari detik pertama ke detik satu seperseratus. --- Hampir dua.

--- Dari detik pertama ke detik pertama. --- Tidak bisa. Penyebutnya nol.

Berhenti di situ. Yang berubah bukan jawabannya, melainkan jarak antara dua titik yang
dibandingkan. Jarak mengecil, garis potong berputar, angka menetap: nilai yang tidak pernah
dicapai tetapi selalu didekati, itulah turunan.

Ulangi bersama: jarak mengecil, garis berputar, angka menetap.
\end{bacakeras}

Mari hitung satu turunan dari definisi. Untuk $f(x)=x^{2}$ di $c=2$,
\[
  \frac{f(2+h)-f(2)}{h}=\frac{(4+4h+h^{2})-4}{h}=\frac{h(4+h)}{h}=4+h .
\]
Karena $h\neq0$, pembagian itu sah, dan limitnya $4$. Untuk $h=0{,}01$ hasil baginya
$4{,}01$; selisih $0{,}01$ itu bukan kesalahan hitung, melainkan suku orde $h$ yang hilang
hanya di limit. Untuk $f(x)=1/x$ di $c=2$,
\[
  \frac{1/(2+h)-1/2}{h}
  =\frac{2-(2+h)}{2h(2+h)}=\frac{-1}{2(2+h)}\longrightarrow-\frac14 .
\]
Pada kedua hitungan aljabar menghapus faktor $h$ dari penyebut, sehingga bentuk $0/0$
berubah menjadi fungsi yang kontinu di $h=0$.

Bentuk hasil bagi itu penting, tetapi tidak bertahan ketika variabelnya menjadi banyak.
Bentuk yang bertahan adalah aproksimasi linear.

\begin{definisi}[Turunan sebagai koefisien aproksimasi linear]
Fungsi $f$ terdiferensialkan di $c$ dengan turunan $L$ bila
\[
  f(c+h)=f(c)+Lh+r(h),
  \qquad\text{dengan}\qquad
  \lim_{h\to0}\frac{r(h)}{h}=0 .
\]
\end{definisi}

Sisa $r(h)$ lebih kecil daripada $h$; itulah isi definisinya. Untuk satu variabel bentuk ini
terasa berlebihan, tetapi hanya bentuk inilah yang bertahan sampai bagian Jacobian.

\begin{proposisi}
Kedua definisi turunan di atas setara, dan bilangan $L$ pada definisi kedua tunggal.
\end{proposisi}

\begin{proof}
Bila limit hasil bagi sama dengan $L$, tetapkan $r(h)=f(c+h)-f(c)-Lh$. Maka $r(h)/h$ tepat
sama dengan hasil bagi dikurangi $L$, sehingga menuju $0$. Sebaliknya, bila $r$ dengan sifat
itu ada, hasil baginya sama dengan $L+r(h)/h$, yang menuju $L$.

Untuk ketunggalan, misalkan $L_{1}$ dan $L_{2}$ sama-sama memenuhi. Maka
$(L_{1}-L_{2})h=r_{2}(h)-r_{1}(h)$, dan setelah dibagi $h$ kita peroleh
$L_{1}-L_{2}=r_{2}(h)/h-r_{1}(h)/h$. Ruas kanan menuju $0$, sedangkan ruas kiri tidak
bergantung pada $h$. Jadi $L_{1}=L_{2}$.
\end{proof}

Sifat pertama yang perlu dipegang: diferensiabilitas lebih kuat daripada kekontinuan.

\begin{proposisi}
Bila $f$ terdiferensialkan di $c$, maka $f$ kontinu di $c$.
\end{proposisi}

\begin{proof}
Untuk $h\neq0$ berlaku
\[
  f(c+h)-f(c)=h\cdot\frac{f(c+h)-f(c)}{h}.
\]
Faktor kedua menuju $f'(c)$ dan faktor pertama menuju $0$, jadi hasil kalinya menuju $0$.
Dengan kata lain $\lim_{h\to0}f(c+h)=f(c)$.
\end{proof}

Kebalikannya tidak berlaku. Fungsi $|x|$ kontinu di $0$ tetapi tidak terdiferensialkan di
sana, karena hasil bagi duanya berbeda tanda. Fungsi $x\sin(1/x)$, bernilai $0$ di $x=0$,
juga kontinu tetapi tidak terdiferensialkan di sana: hasil baginya $\sin(1/x)$, yang berayun
tanpa menuju satu nilai.

\begin{tantangan}
Misalkan $f(0)=0$ dan, untuk $x\neq0$,
\[
  f(x)=|x|\bigl(2+\sin(1/x)\bigr).
\]
Tunjukkan bahwa $f$ kontinu di $0$ tetapi tidak punya turunan di sana. Kedua sifat tidak
bertentangan: kekontinuan menuntut nilai fungsi mendekati $f(0)$, turunan menuntut hasil
bagi selisih mendekati satu bilangan.

Petunjuk 1: untuk kekontinuan, cukup satu ketaksamaan yang mengapit $|f(x)|$ di antara $0$
dan sesuatu yang menuju $0$.

Petunjuk 2: untuk yang kedua, hitung $f(h)/h$ terpisah bagi $h>0$ dan $h<0$. Pada $h>0$
hasil baginya $2+\sin(1/h)$, yang berayun antara $1$ dan $3$; pada $h<0$ ia berayun antara
$-3$ dan $-1$. Kedua kumpulan nilai itu tidak beririsan.

Kalau mentok, kembali ke definisi turunan di awal bagian ini dan hitung hasil baginya untuk
$h>0$ lalu $h<0$ pada fungsi $|x|$.
\end{tantangan}

Aturan hitung berikut membebaskan kita dari definisi setiap kali.

\begin{teorema}[Aturan dasar]
Misalkan $f$ dan $g$ terdiferensialkan di $c$. Maka
\begin{enumerate}
  \item $(f+g)'(c)=f'(c)+g'(c)$ dan $(af)'(c)=a f'(c)$ untuk $a\in\mathbb{R}$;
  \item $(fg)'(c)=f'(c)g(c)+f(c)g'(c)$;
  \item $\left(\dfrac{f}{g}\right)'(c)=\dfrac{g(c)f'(c)-f(c)g'(c)}{g(c)^{2}}$ bila
        $g(c)\neq0$;
  \item bila $n$ bilangan bulat, turunan $x^{n}$ di $x\neq0$ adalah $nx^{n-1}$.
\end{enumerate}
\end{teorema}

\begin{proof}
Bagian (1) langsung dari limit jumlah. Untuk (2), tulis
\[
  f(x)g(x)-f(c)g(c)=f(x)\bigl[g(x)-g(c)\bigr]+g(c)\bigl[f(x)-f(c)\bigr].
\]
Setelah dibagi $x-c$ dan $x\to c$, faktor $f(x)$ menuju $f(c)$ karena $f$ kontinu di $c$
(proposisi sebelumnya), dan kedua hasil bagi menuju turunannya.

Untuk (3), cukup ditunjukkan bahwa $1/g$ terdiferensialkan dengan turunan $-g'(c)/g(c)^{2}$,
lalu pakai (2) pada $f\cdot(1/g)$. Hasil bagi untuk $1/g$ adalah
\[
  \frac{1/g(x)-1/g(c)}{x-c}
  =-\frac{1}{g(x)g(c)}\cdot\frac{g(x)-g(c)}{x-c},
\]
yang menuju $-g'(c)/g(c)^{2}$ karena $g$ kontinu dan $g(c)\neq0$. Bagian (4) diperoleh dari
(2) dengan induksi pada $n\ge0$, lalu dari (3) untuk $n<0$.
\end{proof}

\section{Aturan rantai}
\tingkat{menengah}

Bila satu besaran bergantung pada besaran kedua, dan besaran kedua bergantung pada waktu,
maka laju perubahan besaran pertama adalah hasil kali dua laju itu. Kalimat itu terdengar
seperti aritmetika sehari-hari, dan memang itulah isi aturan rantai.

Gambar~\ref{fig:roda-gigi} memakai dua roda gigi: roda dalam berputar tiga kali per putaran
roda tengah, dan roda tengah dua kali per putaran roda luar. Satu putaran roda luar memutar
roda dalam $3\times2=6$ kali; menjumlahkan kedua rasio memberi $5$, angka yang tidak muncul
pada putaran mana pun.

\begin{figure}[htbp]
\centering
\begin{tikzpicture}[scale=0.95]
% roda dalam, rasio 3
\draw[ink, thick] (0,0) circle (1.05);
\foreach \a in {0,20,...,340} { \draw[ink, thick] (\a:1.05) -- (\a:1.17); }
\fill[ink] (0,0) circle (1.5pt);
\draw[->, layerindigo, thick] (0,0) ++(150:0.55) arc (150:70:0.55);
% roda luar, rasio 2
\begin{scope}[xshift=2.9cm]
  \draw[ink, thick] (0,0) circle (1.8);
  \foreach \a in {10,30,...,350} { \draw[ink, thick] (\a:1.8) -- (\a:1.95); }
  \fill[ink] (0,0) circle (1.5pt);
  \draw[->, layerbrown, thick] (0,0) ++(-20:0.75) arc (-20:-100:0.75);
\end{scope}
\node[font=\footnotesize, text=ink, align=center] at (0,-2.35) {roda dalam:\\3 putaran};
\node[font=\footnotesize, text=ink, align=center] at (2.9,-2.35) {roda luar:\\2 putaran};
\node[font=\footnotesize, text=layerplum] at (1.45,2.15) {$3\times2=6$};
\end{tikzpicture}
\caption{Aturan rantai sebagai dua roda gigi. Roda dalam berputar $3$ kali untuk satu putaran
roda luar, dan roda luar $2$ kali untuk satu putaran keluarannya; rasio totalnya
$3\times2=6$, bukan $3+2=5$.}
\label{fig:roda-gigi}
\end{figure}

\kamus{aturan rantai}{chain rule}{Laju di ujung rantai sama dengan hasil kali laju pada
setiap mata rantai.}

\begin{teorema}[Aturan rantai]
Misalkan $g$ terdiferensialkan di $c$ dan $f$ terdiferensialkan di $g(c)$. Maka $f\circ g$
terdiferensialkan di $c$ dan
\[
  (f\circ g)'(c)=f'\bigl(g(c)\bigr)\,g'(c).
\]
\end{teorema}

\begin{proof}
Tulis $y=g(c)$. Karena $f$ terdiferensialkan di $y$, menurut definisi aproksimasi linear
ada fungsi $v$ dengan $v(s)\to0$ untuk $s\to y$ dan
\[
  f(s)-f(y)=(s-y)\bigl[f'(y)+v(s)\bigr].
\]
Dengan cara yang sama, karena $g$ terdiferensialkan di $c$, ada $u$ dengan $u(t)\to0$ untuk
$t\to c$ dan
\[
  g(t)-g(c)=(t-c)\bigl[g'(c)+u(t)\bigr].
\]
Substitusikan $s=g(t)$ pada persamaan pertama, lalu pakai persamaan kedua:
\[
  f\bigl(g(t)\bigr)-f\bigl(g(c)\bigr)
  =(t-c)\bigl[g'(c)+u(t)\bigr]\bigl[f'(y)+v(g(t))\bigr].
\]
Untuk $t\neq c$, bagi dengan $t-c$. Karena $g$ kontinu di $c$, $g(t)\to y$, sehingga
$v(g(t))\to0$, dan hasil baginya menuju $g'(c)f'(y)$.
\end{proof}

Pembuktian naif yang membagi dengan $g(t)-g(c)$ gagal ketika $g(t)=g(c)$ untuk $t$ dekat
$c$ tetapi $t\neq c$; fungsi konstan lokal, misalnya, membuat penyebut itu nol di semua
titik. Penyangga $v$ di atas menghindari pembagian tersebut sama sekali.

Contoh pendek: untuk $h(x)=(x^{2}+1)^{5}$, ambil $g(x)=x^{2}+1$ dan $f(u)=u^{5}$. Maka
$h'(x)=5(x^{2}+1)^{4}\cdot2x$ dan $h'(1)=5\cdot16\cdot2=160$; ekspansi langsung mungkin,
tetapi tidak menyenangkan.

\begin{resep}
\begin{enumerate}
  \item \textbf{Pisahkan dua lapis.} Tunjuk mana yang di dalam dan mana yang di luar. Untuk
        $k(x)=\sin(x^{3}+2x)$: lapis dalam $g(x)=x^{3}+2x$, lapis luar $f(u)=\sin u$.
  \item \textbf{Turunkan lapis luar, biarkan isinya utuh.} $f'(u)=\cos u$ menjadi
        $\cos(x^{3}+2x)$; jangan dulu menyederhanakan isinya.
  \item \textbf{Kalikan turunan lapis dalam.} $g'(x)=3x^{2}+2$, sehingga
        $k'(x)=(3x^{2}+2)\cos(x^{3}+2x)$.
  \item \textbf{Uji satu titik.} Di $x=0$ rumus memberi $2$, dan beda pusat dengan
        $h=10^{-3}$ memberi $1{,}9999997$.
\end{enumerate}
Peringatan urutan: lapis luar dinilai di $g(x)$, bukan di $x$. Menukar tempat penilaian
memberi $\cos(x)\,(3x^{2}+2)$, yang di $x=1$ bernilai $2{,}7015$, sedangkan nilai benarnya
$-4{,}9499$.
\end{resep}

Contoh laju: luas lingkaran $A=\pi r^{2}$ dengan $\dot r=0{,}5$ cm per detik. Di $r=10$ cm,
\[
  \frac{dA}{dt}=2\pi r\frac{dr}{dt}=2\pi\cdot10\cdot0{,}5=10\pi\approx31{,}4159
  \ \text{cm}^{2}\ \text{per detik}.
\]
Angka $10\pi$ eksak; di $r=20$ cm laju yang sama memberi $20\pi$: jawabannya bergantung
pada $r$.

\section{Teorema nilai rata-rata dan akibatnya}
\tingkat{menengah}

Di antara dua titik pada grafik yang mulus selalu ada satu titik yang kemiringannya sama
dengan kemiringan talibusur yang menghubungkan kedua titik itu. Pernyataan itu terlihat
seperti khayalan geometris, tetapi dialah jembatan antara turunan dan fungsi.

Kunci pembuktiannya adalah kasus khusus ketika kedua ujung bernilai sama.

\begin{teorema}[Rolle]
Bila $f$ kontinu pada $[a,b]$, terdiferensialkan pada $(a,b)$, dan $f(a)=f(b)$, maka ada
$c\in(a,b)$ dengan $f'(c)=0$.\index{teorema nilai rata-rata}
\end{teorema}

\begin{proof}
Bila $f$ konstan, setiap $c\in(a,b)$ memenuhi. Bila tidak, $f$ mencapai nilai maksimum atau
minimum yang berbeda dari nilai ujungnya; karena $f$ kontinu pada interval tertutup terbatas,
nilai itu dicapai di suatu titik, dan titik itu tidak boleh ujung, jadi $c\in(a,b)$ (teorema
nilai ekstrem, Bab 11). Di titik ekstrem interior, hasil bagi
$(f(c+h)-f(c))/h$ bernilai $\le0$ untuk $h>0$ kecil dan $\ge0$ untuk $h<0$ kecil. Kedua
ketaksamaan itu bersama memaksa limit duanya, yaitu $f'(c)$, bernilai $0$.
\end{proof}

\begin{teorema}[Nilai rata-rata]
Bila $f$ kontinu pada $[a,b]$ dan terdiferensialkan pada $(a,b)$, maka ada $c\in(a,b)$
dengan
\[
  f'(c)=\frac{f(b)-f(a)}{b-a}.
\]
\end{teorema}

\begin{proof}
Tetapkan $s=(f(b)-f(a))/(b-a)$ dan $g(x)=f(x)-f(a)-s(x-a)$. Maka $g$ kontinu pada $[a,b]$,
terdiferensialkan pada $(a,b)$, dan $g(a)=g(b)=0$. Menurut teorema Rolle ada $c\in(a,b)$
dengan $g'(c)=0$, yaitu $f'(c)=s$.
\end{proof}

Contoh angka: $f(x)=x^{3}$ pada $[0,3]$ berkemiringan talibusur $27/3=9$, dan $3c^{2}=9$
memberi $c=\sqrt{3}\approx1{,}7321$; tidak ada jaminan $c$ di tengah.

\begin{tangan}
Siapkan papan tipis 30 cm, dua buku, penggaris, dan pensil. Ganjal ujung kiri papan setinggi
4 cm, lalu tekan sedikit bagian tengahnya supaya melengkung ke atas.

Tegakkan penggaris pada 5 cm dan 25 cm dari ujung yang diangkat; selisih tinggi dibagi 20 cm
memberi satu angka kemiringan. Ulangi dengan 5 cm dan 10 cm: angkanya berubah, sebab yang
diukur adalah kemiringan rata-rata pada selang berbeda.

Sekarang letakkan penggaris rata di titik 15 cm dan putar sampai hanya satu titik yang
menempel. Angkat ujung yang lain perlahan: penggaris berputar pada titik itu, persis seperti
talibusur pada Gambar~\ref{fig:sekan-singgung}. Kemiringan penggaris yang menyentuh satu
titik itulah kemiringan di titik itu, dan di situlah teorema nilai rata-rata menjanjikan
adanya satu titik semacam itu.
\end{tangan}

\begin{akibat}[Uji monotonisitas]
Misalkan $f$ kontinu pada $[a,b]$ dan terdiferensialkan pada $(a,b)$. Bila $f'(x)\ge0$ untuk
semua $x\in(a,b)$, maka $f$ tidak turun pada $[a,b]$. Bila $f'(x)>0$ untuk semua $x$, maka
$f$ naik tegas. Bila $f'(x)=0$ untuk semua $x$, maka $f$ konstan.
\end{akibat}

\begin{proof}
Ambil $a\le x<y\le b$. Teorema nilai rata-rata pada $[x,y]$ memberi $c\in(x,y)$ dengan
$f(y)-f(x)=f'(c)(y-x)$. Karena $y-x>0$, tanda selisih nilai mengikuti tanda $f'(c)$. Bila
$f'$ nol di mana-mana, selisihnya nol untuk setiap pasangan, jadi $f$ konstan.
\end{proof}

Arah logikanya perlu diperhatikan: akibat ini mengubah informasi tentang turunan menjadi
informasi tentang fungsi. Kebalikannya tidak berlaku: fungsi yang naik tegas boleh punya
turunan nol di beberapa titik, seperti $x^{3}$ di $0$. Juga tidak berlaku kesimpulan balik
dari satu titik: $f'(c)=0$ tidak berarti $c$ titik ekstrem.

\begin{akibat}[Keterkendalian]
Bila $f$ terdiferensialkan pada interval $I$ dan $|f'(x)|\le M$ untuk semua $x\in I$, maka
$|f(x)-f(y)|\le M|x-y|$ untuk semua $x,y\in I$.
\end{akibat}

\begin{proof}
Untuk $x\neq y$, teorema nilai rata-rata pada interval di antara keduanya memberi $c$
dengan $f(x)-f(y)=f'(c)(x-y)$. Ambil nilai mutlak kedua ruas dan pakai $|f'(c)|\le M$.
\end{proof}

Karena itu fungsi dengan turunan terbatas otomatis seragam kontinu (ambil
$\delta=\varepsilon/M$), tetapi kebalikannya tidak berlaku: $\sqrt{x}$ seragam kontinu pada
$[0,1]$ dengan turunan tak terbatas di dekat nol.

\section{Polinomial Taylor dan sisanya}
\tingkat{lanjut}

Garis singgung adalah polinomial berderajat satu yang menempel pada grafik di satu titik.
Bila kita mau lebih lama menempel, kita tambah suku berderajat dua, tiga, dan seterusnya.
Pertanyaannya bukan apakah penambahan itu memperbaiki hampiran, tetapi seberapa besar
sisanya.

\begin{definisi}[Polinomial Taylor]
Misalkan $f$ mempunyai turunan sampai orde $n$ di $a$. Polinomial Taylor orde $n$ dari $f$
di $a$ adalah
\[
  P_{n}(x)=\sum_{k=0}^{n}\frac{f^{(k)}(a)}{k!}(x-a)^{k}.
\]
\end{definisi}

Syaratnya: $P_{n}$ dan $f$ sama nilai dan $n$ turunan pertamanya di $a$; karena $n+1$
koefisien menentukan satu polinomial berderajat $n$, $P_{n}$ tunggal.

\begin{teorema}[Taylor dengan sisa Lagrange]
Misalkan $f$ mempunyai turunan sampai orde $n-1$ yang kontinu pada $[a,b]$ dan turunan
ke-$n$ pada $(a,b)$. Untuk setiap $\beta\in(a,b]$ ada $\xi\in(a,\beta)$ dengan
\[
  f(\beta)=P_{n-1}(\beta)+\frac{f^{(n)}(\xi)}{n!}(\beta-a)^{n}.
\]
\end{teorema}

\begin{proof}
Tetapkan $M$ oleh $f(\beta)=P_{n-1}(\beta)+M(\beta-a)^{n}$, lalu
\[
  g(t)=f(t)-P_{n-1}(t)-M(t-a)^{n},\qquad a\le t\le b .
\]
$P_{n-1}$ dibangun dari turunan $f$ di $a$, jadi $g(a)=\cdots=g^{(n-1)}(a)=0$, dan pilihan
$M$ memberi $g(\beta)=0$. Rolle memberi $x_{1}\in(a,\beta)$ dengan $g'(x_{1})=0$; karena
$g'(a)=0$, Rolle pada $g'$ di $[a,x_{1}]$ memberi $x_{2}$ dengan $g''(x_{2})=0$. Ulangi sampai
$n$ langkah: ada $x_{n}\in(a,\beta)$ dengan $g^{(n)}(x_{n})=0$. Turunan ke-$n$ dari $P_{n-1}$
nol dan dari $(t-a)^{n}$ adalah $n!$, sehingga $0=f^{(n)}(x_{n})-n!M$, yaitu
$n!M=f^{(n)}(\xi)$ dengan $\xi=x_{n}$.
\end{proof}

Untuk $n=1$ teorema ini tepat teorema nilai rata-rata; yang baru adalah kemampuan menghitung
galat orde tinggi.

\begin{contoh}
Hampiri $e$ dengan $P_{3}(1)$ dari $f(x)=e^{x}$ di $a=0$, lalu bandingkan dengan batas
sisanya.
\end{contoh}

Di sini $f^{(k)}(0)=1$ untuk semua $k$, sehingga
\[
  P_{3}(1)=1+1+\frac{1}{2}+\frac{1}{6}=\frac{8}{3}=2{,}6666\ldots
\]
Nilai sebenarnya $e=2{,}71828\ldots$, jadi galatnya $0{,}0516152$. Sisa Lagrange memberi
$|e-P_{3}(1)|=e^{\xi}/24$ untuk suatu $\xi\in(0,1)$, dan karena $e^{\xi}\le e<3$, batasnya
$3/24=0{,}125$. Batas itu benar, tetapi longgar sekitar dua setengah kali. Memakai orde
lima, $P_{5}(1)=2{,}71666\ldots$ dengan galat $0{,}0016152$, sedangkan batasnya
$3/720=0{,}0041667$.

\begin{untukanak}
Bayangkan kamu menebak jalan kawanmu hanya dari beberapa petunjuk. Petunjuk pertama: di mana
dia sekarang. Tebakanmu satu titik, dan titik itu tidak bergerak. Petunjuk kedua: ke arah mana
dia melangkah. Tebakanmu garis lurus, sudah lebih mirip, tetapi jalan sungguhan biasanya
melengkung. Petunjuk ketiga: seberapa cepat arahnya berubah. Tebakanmu potongan parabola, yang
menempel lebih lama.

Semakin banyak petunjuk, semakin lama tebakanmu menempel, tetapi tidak pernah sama persis.
Selalu ada sisa. Sisa itu dapat dibatasi tanpa mengetahui jalan sebenarnya, hanya dari
seberapa besar perubahan pada tingkat berikutnya. Di bab ini kita akan tahu mengapa garis
lurus tebakan orde satu, mengapa parabola orde dua, dan mengapa batas sisa lebih berguna
daripada jawaban yang katanya tepat.
\end{untukanak}

\section{Turunan parsial, gradien, dan turunan berarah}
\tingkat{menengah}

Bila masukan lebih dari satu, ``laju perubahan'' belum tentu satu bilangan. Ia bergantung
pada arah kita melangkah. Cara paling sederhana mendapat satu bilangan adalah membekukan
semua koordinat kecuali satu, lalu memakai definisi yang sudah kita punya.

\begin{definisi}[Turunan parsial]
Misalkan $f\colon U\to\mathbb{R}$ terdefinisi pada himpunan terbuka
$U\subseteq\mathbb{R}^{n}$ dan $a\in U$. Turunan parsial $f$ terhadap $x_{j}$ di $a$ adalah
\[
  \frac{\partial f}{\partial x_{j}}(a)
  =\lim_{h\to0}\frac{f(a+h\,e_{j})-f(a)}{h},
\]
dengan $e_{j}$ vektor satuan arah $x_{j}$, asalkan limitnya ada.\index{turunan parsial}
\end{definisi}

Secara geometris, turunan parsial terhadap $x$ adalah kemiringan irisan grafik oleh bidang
$y$ tetap, jadi turunan biasa fungsi satu variabel.

\kamus{gradien}{gradient}{Panah yang menunjuk arah paling curam pada sebuah bukit, dengan
panjang sebesar kecuraman di arah itu.}

\begin{definisi}[Gradien]
Bila semua turunan parsial $f$ di $a$ ada, vektor
\[
  \nabla f(a)=\left(\frac{\partial f}{\partial x_{1}}(a),\dots,
  \frac{\partial f}{\partial x_{n}}(a)\right)
\]
disebut gradien $f$ di $a$.\index{gradien}
\end{definisi}

Untuk $f(x,y)=x^{2}y^{3}$ berlaku $\partial f/\partial x=2xy^{3}$ dan
$\partial f/\partial y=3x^{2}y^{2}$, sehingga $\nabla f(2,-1)=(-4,12)$.

\begin{definisi}[Turunan berarah]
Untuk vektor satuan $u\in\mathbb{R}^{n}$, turunan berarah $f$ di $a$ sepanjang $u$ adalah
\[
  D_{u}f(a)=\lim_{t\to0}\frac{f(a+tu)-f(a)}{t},
\]
asalkan limitnya ada.\index{turunan berarah}
\end{definisi}

Turunan parsial adalah kasus khusus $u=e_{j}$. Untuk arah lain $D_{u}f$ tidak dapat diduga
dari daftar parsial, kecuali dalam satu keadaan baik.

\begin{proposisi}
Bila $f$ terdiferensialkan di $a$ (dalam pengertian aproksimasi linear di bagian berikut),
maka untuk setiap vektor satuan $u$ berlaku $D_{u}f(a)=\nabla f(a)\cdot u$.
\end{proposisi}

\begin{proof}
Dengan $v=tu$ pada definisi aproksimasi linear,
$f(a+tu)=f(a)+t\,\nabla f(a)\cdot u+r(tu)$, sehingga
\[
  \frac{f(a+tu)-f(a)}{t}=\nabla f(a)\cdot u+\frac{r(tu)}{t}.
\]
Karena $\|r(tu)\|/\|tu\|\to0$ dan $\|tu\|=|t|$, suku terakhir menuju $0$.
\end{proof}

\begin{proposisi}[Arah naik tercepat]
Bila $\nabla f(a)\neq0$, maka di antara semua vektor satuan $u$, nilai $D_{u}f(a)$ terbesar
dicapai pada $u=\nabla f(a)/\|\nabla f(a)\|$, dan nilai terbesarnya $\|\nabla f(a)\|$.
\end{proposisi}

\begin{proof}
Menurut ketaksamaan Cauchy--Schwarz (Bab 10),
\[
  |D_{u}f(a)|=|\nabla f(a)\cdot u|\le\|\nabla f(a)\|\,\|u\|=\|\nabla f(a)\|.
\]
Kesamaan berlaku tepat ketika $u$ sejajar dengan $\nabla f(a)$. Karena $u$ vektor satuan,
hanya dua tanda yang mungkin pada arah itu, dan tanda positif memberi nilai terbesar.
\end{proof}

Untuk $f(x,y)=x^{2}y^{3}$ di $(2,-1)$, panjang gradiennya $\sqrt{160}\approx12{,}6491$, jadi
arah naik tercepat $(-4,12)/\sqrt{160}\approx(-0{,}3162,\,0{,}9487)$ dengan laju $12{,}6491$
per satuan jarak. Arah $x$ memberi $-4$, arah $y$ memberi $12$; keduanya lebih kecil.

Bila kurva $t\mapsto\gamma(t)$ tinggal pada himpunan nivel $f=\text{tetap}$, maka
$f(\gamma(t))$ konstan dan turunannya nol; aturan rantai memberi
$\nabla f(\gamma(t))\cdot\gamma'(t)=0$. Jadi gradien tegak lurus garis nivel, dan arah naik
tercepat tegak lurus garis kontur pada peta.

Contohnya $f(x,y)=x^{2}+2y^{2}$ di $(2,-1)$, titik dengan $f=6$, sehingga ia duduk pada
elips nivel $x^{2}+2y^{2}=6$. Gradiennya $(2x,4y)=(4,-4)$, panjangnya
$4\sqrt{2}\approx5{,}6569$, sedangkan arah singgung elips di sana $(1,1)$ dan
$(4,-4)\cdot(1,1)=0$. Gambar~\ref{fig:kontur-gradien} menggambar ketiganya: elips nivel,
panah gradien, dan garis singgung yang tegak lurus padanya. Panah gradien tidak menyusuri
kontur; ia menembusnya ke arah nilai $f$ bertambah.

\begin{figure}[htbp]
\centering
\begin{tikzpicture}[scale=0.92]
\draw[->, layergrey] (-3.0,0) -- (3.4,0) node[right, font=\scriptsize] {$x$};
\draw[->, layergrey] (0,-2.45) -- (0,2.5) node[above, font=\scriptsize] {$y$};
\draw[ink] (0,0) ellipse [x radius=1.414, y radius=1.000];
\draw[ink] (0,0) ellipse [x radius=2.000, y radius=1.414];
\draw[ink, thick] (0,0) ellipse [x radius=2.449, y radius=1.732];
\node[font=\scriptsize, anchor=west, text=ink] at (0.28,-0.98) {$c=2$};
\node[font=\scriptsize, anchor=west, text=ink] at (0.28,-1.40) {$c=4$};
\node[font=\scriptsize, anchor=west, text=ink] at (0.28,-1.98) {$c=6$};
\draw[accent, thick, dashed] (1.30,-1.70) -- (2.72,-0.28);
\draw[->, very thick, layerindigo] (2,-1) -- (2.707,-1.707);
\fill[ink] (2,-1) circle (1.7pt);
\node[font=\scriptsize, anchor=north west, text=ink] at (2.02,-1.05) {$(2,-1)$};
\node[font=\scriptsize, anchor=west, text=layerindigo] at (2.72,-1.86) {$\nabla f=(4,-4)$};
\node[font=\scriptsize, anchor=east, text=accent] at (1.32,-1.76) {arah $(1,1)$};
\end{tikzpicture}
\caption{Kurva nivel $x^{2}+2y^{2}=c$ untuk $c=2,4,6$ dan gradien di $(2,-1)$. Panah
gradien $(4,-4)$ tegak lurus arah singgung $(1,1)$ pada kurva nivel $c=6$, sebab
$(4,-4)\cdot(1,1)=0$.}
\label{fig:kontur-gradien}
\end{figure}

\section{Diferensiabilitas, Jacobian, dan aturan rantai multivariabel}
\tingkat{lanjut}

Turunan yang sesungguhnya bukan daftar parsial, melainkan satu peta linear yang meniru fungsi
di sekitar titik; daftar parsial dapat menyesatkan.

\begin{definisi}[Diferensiabilitas]
Misalkan $U\subseteq\mathbb{R}^{n}$ terbuka dan $F\colon U\to\mathbb{R}^{m}$. $F$
terdiferensialkan di $a\in U$ bila ada peta linear $L\colon\mathbb{R}^{n}\to\mathbb{R}^{m}$
dengan
\[
  F(a+v)=F(a)+Lv+r(v),
  \qquad
  \lim_{v\to0}\frac{\|r(v)\|}{\|v\|}=0 .
\]
Peta $L$ disebut turunan $F$ di $a$ dan ditulis $F'(a)$.\index{diferensiabilitas}
\end{definisi}

\begin{proposisi}
Peta $L$ pada definisi itu tunggal bila ada.
\end{proposisi}

\begin{proof}
Bila $L_{1}$ dan $L_{2}$ memenuhi, maka $(L_{1}-L_{2})v=r_{2}(v)-r_{1}(v)$. Ambil
$v=tw$ dengan $w$ tetap dan $t>0$, bagi dengan $t$, lalu $t\to0$: diperoleh
$(L_{1}-L_{2})w=0$. Karena berlaku untuk setiap $w$, kedua peta sama.
\end{proof}

\begin{teorema}
Bila $F$ terdiferensialkan di $a$, maka semua turunan parsialnya di $a$ ada, dan matriks
$L$ terhadap basis standar adalah
\[
  J_{F}(a)=\begin{pmatrix}
    \dfrac{\partial F_{1}}{\partial x_{1}}(a) & \cdots &
    \dfrac{\partial F_{1}}{\partial x_{n}}(a)\\
    \vdots & \ddots & \vdots\\
    \dfrac{\partial F_{m}}{\partial x_{1}}(a) & \cdots &
    \dfrac{\partial F_{m}}{\partial x_{n}}(a)
  \end{pmatrix}.
\]
\end{teorema}

\begin{proof}
Ambil $v=te_{j}$ dengan $t\neq0$. Definisi memberi
$F(a+te_{j})-F(a)=tLe_{j}+r(te_{j})$, dan setelah dibagi $t$ lalu $t\to0$, kolom ke-$j$
dari $L$ muncul sebagai limit hasil bagi, yaitu vektor turunan parsial terhadap $x_{j}$.
\end{proof}

\kamus{Jacobian}{Jacobian matrix}{Tabel semua kemiringan sebuah fungsi banyak variabel:
setiap baris satu keluaran, setiap kolom satu masukan.}

Matriks $m\times n$ itu disebut matriks Jacobian $F$ di $a$.\index{jacobian} Untuk $m=1$ ia
menjadi baris tunggal, yaitu transpos gradien. Contoh: $F(x,y)=(x^{2}y,\;x+y^{2})$ di $(1,2)$
memberi
\[
  J_{F}(1,2)=\begin{pmatrix}2xy & x^{2}\\ 1 & 2y\end{pmatrix}_{(1,2)}
  =\begin{pmatrix}4&1\\ 1&4\end{pmatrix},
\]
sehingga $F(1{,}01,\,1{,}98)\approx(2,5)+J_{F}(1,2)(0{,}01,-0{,}02)=(2{,}02,\,4{,}93)$,
sedangkan nilai eksaknya $(2{,}019798,\,4{,}9304)$; hampiran itu meleset $2\times10^{-4}$
pada komponen pertama.

Parsial yang ada tidak menjamin diferensiabilitas; dua contoh berikut menunjukkan dua tingkat
kegagalan.

\begin{contoh}[Parsial ada, tidak kontinu]
Misalkan $f(x,y)=\dfrac{xy}{x^{2}+y^{2}}$ untuk $(x,y)\neq(0,0)$ dan $f(0,0)=0$. Di titik
asal $f(h,0)=0$, sehingga kedua turunan parsialnya nol. Tetapi sepanjang garis $y=kx$ nilai
fungsinya $k/(1+k^{2})$: untuk $k=1$ nilainya $1/2$, untuk $k=0$ nilainya $0$. Limitnya tidak
ada, dan $f$ tidak kontinu di titik asal.
\end{contoh}

\begin{contoh}[Kontinu, parsial ada, tak terdiferensialkan]
Misalkan $g(x,y)=\dfrac{xy}{\sqrt{x^{2}+y^{2}}}$ untuk $(x,y)\neq(0,0)$ dan $g(0,0)=0$.
Karena $2|xy|\le x^{2}+y^{2}$, berlaku $|g(x,y)|\le\frac12\sqrt{x^{2}+y^{2}}$, jadi $g$
kontinu di titik asal, dan kedua parsialnya di sana nol. Bila $g$ terdiferensialkan, peta
linearnya harus nol, sehingga sisa $r(v)=g(v)$ mesti memenuhi $\|r(v)\|/\|v\|\to0$. Ambil
$v=t(1,1)/\sqrt{2}$ dengan $t>0$: maka $g(v)=(t^{2}/2)/t=t/2$, sehingga
$g(v)/\|v\|=1/2$ untuk setiap $t$. Nisbah itu menuju $1/2$, bukan $0$; jadi $g$ tidak
terdiferensialkan di titik asal.
\end{contoh}

Akibatnya penting: pada contoh kedua $\nabla g(0,0)=0$ tetapi turunan berarah diagonal tidak
ada, karena hasil baginya $\pm1/2$ bergantung tanda $t$. Jadi $D_{u}f=\nabla f\cdot u$
mensyaratkan diferensiabilitas.

Syarat yang mudah diperiksa dan cukup untuk diferensiabilitas adalah kekontinuan parsial.

\begin{teorema}
Bila turunan-turunan parsial $F$ ada dan kontinu di sekitar $a$, maka $F$ terdiferensialkan
di $a$.
\end{teorema}

\begin{proof}
Cukup dibuktikan untuk $m=1$ dan $n=2$; kasus lain mengikuti komponen demi komponen dengan
indeks yang bertambah. Tulis $a=(p,q)$ dan $v=(h,k)$. Maka
\[
  f(p+h,q+k)-f(p,q)
  =\bigl[f(p+h,q+k)-f(p,q+k)\bigr]+\bigl[f(p,q+k)-f(p,q)\bigr].
\]
Menurut teorema nilai rata-rata satu variabel, suku pertama sama dengan
$h\,\partial_{x}f(\xi,q+k)$ untuk suatu $\xi$ di antara $p$ dan $p+h$, dan suku kedua sama
dengan $k\,\partial_{y}f(p,\eta)$ untuk suatu $\eta$ di antara $q$ dan $q+k$. Karena itu
\[
  f(a+v)-f(a)-\nabla f(a)\cdot v
  =h\bigl[\partial_{x}f(\xi,q+k)-\partial_{x}f(a)\bigr]
  +k\bigl[\partial_{y}f(p,\eta)-\partial_{y}f(a)\bigr].
\]
Bagi dengan $\|v\|=\sqrt{h^{2}+k^{2}}$: karena $|h|/\|v\|\le1$ dan $|k|/\|v\|\le1$, dan
kedua selisih dalam kurung menuju $0$ (di sinilah kekontinuan parsial dipakai), ruas kanan
dibagi $\|v\|$ menuju $0$.
\end{proof}

Sekarang komposisi.

\begin{teorema}[Aturan rantai multivariabel]
Misalkan $H\colon\mathbb{R}^{n}\to\mathbb{R}^{k}$ terdiferensialkan di $a$ dan
$G\colon\mathbb{R}^{k}\to\mathbb{R}^{m}$ terdiferensialkan di $H(a)$. Maka $G\circ H$
terdiferensialkan di $a$ dan
\[
  (G\circ H)'(a)=G'(H(a))\,H'(a),
\]
yaitu komposisi dua peta linear; dalam matriks,
$J_{G\circ H}(a)=J_{G}(H(a))\,J_{H}(a)$.
\end{teorema}

\begin{proof}
Tulis $A=H'(a)$, $B=G'(y_{0})$ dengan $y_{0}=H(a)$, dan
\[
  H(a+v)=H(a)+Av+r(v),\qquad G(y_{0}+w)=G(y_{0})+Bw+q(w),
\]
dengan $\|r(v)\|\le\varepsilon(v)\|v\|$ dan $\|q(w)\|\le\eta(w)\|w\|$, di mana
$\varepsilon(v)\to0$ dan $\eta(w)\to0$. Untuk $w=H(a+v)-H(a)=Av+r(v)$ berlaku
$\|w\|\le(\|A\|+\varepsilon(v))\|v\|$, jadi $w\to0$ ketika $v\to0$. Maka
\[
  G\bigl(H(a+v)\bigr)-G(H(a))-BAv
  =B\bigl(H(a+v)-H(a)-Av\bigr)+q(w)
  =B\,r(v)+q(w).
\]
Bagi dengan $\|v\|$:
\[
  \frac{\|B\,r(v)+q(w)\|}{\|v\|}
  \le\|B\|\,\varepsilon(v)+\bigl(\|A\|+\varepsilon(v)\bigr)\eta(w),
\]
dan ruas kanan menuju $0$. Jadi $BA$ memenuhi definisi turunan $G\circ H$ di $a$.
\end{proof}

Bentuk matriksnya menjelaskan mengapa urutan perkalian tidak boleh dibalik: $J_{H}$
berukuran $k\times n$ dan $J_{G}$ berukuran $m\times k$, jadi hanya satu urutan yang dapat
dikalikan. Untuk fungsi skalar hasilnya $\frac{d}{dt}f(\gamma(t))=\nabla f(\gamma(t))\cdot
\gamma'(t)$.

Terakhir, turunan kedua untuk fungsi banyak variabel disimpan dalam satu matriks.

\kamus{Hessian}{Hessian matrix}{Tabel kelengkungan: apakah permukaan melengkung ke atas,
ke bawah, atau ke dua arah sekaligus.}

\begin{definisi}[Matriks Hessian]
Bila semua turunan parsial kedua $f\colon U\to\mathbb{R}$ ada di $a$, matriks
\[
  H_{f}(a)=\left(\frac{\partial^{2}f}{\partial x_{i}\,\partial x_{j}}(a)\right)_{i,j=1}^{n}
\]
disebut matriks Hessian $f$ di $a$.\index{hessian}
\end{definisi}

Untuk fungsi yang parsial keduanya kontinu, $H_{f}(a)$ simetris (teorema Clairaut, dibuktikan
di Stewart pada bagian turunan parsial). Tanpa kekontinuan, kesimpulan itu bisa gagal.

Hessian adalah turunan dari gradien. Bila parsial keduanya kontinu di sekitar titik kritis
$a$ ($\nabla f(a)=0$), maka
\[
  f(a+v)=f(a)+\tfrac12\,v^{\top}H_{f}(a)\,v+o\bigl(\|v\|^{2}\bigr),
\]
sehingga tanda bentuk kuadrat $v\mapsto v^{\top}Hv$ menentukan jenis titik kritis.

\begin{proposisi}[Uji turunan kedua]
Misalkan $f$ mempunyai turunan parsial kedua kontinu di sekitar titik kritis $a$
($\nabla f(a)=0$).
\begin{enumerate}
  \item Bila $H_{f}(a)$ definit positif, $a$ titik minimum lokal tegas.
  \item Bila $H_{f}(a)$ definit negatif, $a$ titik maksimum lokal tegas.
  \item Bila $H_{f}(a)$ mempunyai nilai eigen bertanda berbeda, $a$ titik sadel.
  \item Bila $H_{f}(a)$ semidefinit tetapi singular, uji ini tidak memberi kesimpulan.
\end{enumerate}
\end{proposisi}

Bagian pertama sampai ketiga mengikuti ekspansi kuadratik di atas: pada arah $v^{\top}Hv>0$
fungsi naik, pada arah $v^{\top}Hv<0$ fungsi turun. Bukti lengkapnya ada di Stewart pada
bagian nilai maksimum dan minimum. Bagian keempat benar-benar kosong: $x^{4}+y^{4}$ dan
$x^{4}-y^{4}$ sama-sama punya $H=0$ di titik asal, tetapi yang pertama minimum dan yang kedua
sadel. Untuk dua variabel, $AC-B^{2}>0$ dengan $A>0$ berarti definit positif, $AC-B^{2}<0$
berarti indefinit.

\section{Contoh hitung}
\tingkat{menengah}

Tiga hitungan berikut memakai angka yang sama dengan kode di bagian lab.

Pertama, beda hingga untuk $f(x)=\sin x$ di $x_{0}=1$; nilai eksaknya $\cos 1$, yang
angkanya muncul pada keluaran kode di bagian lab. Dua rumusnya:
\[
  \text{maju: }\frac{f(x_{0}+h)-f(x_{0})}{h},\qquad
  \text{pusat: }\frac{f(x_{0}+h)-f(x_{0}-h)}{2h}.
\]

\begin{center}
\begin{tabular}{rrrrr}
\hline
$h$ & beda maju & galat relatif & beda pusat & galat relatif\\
\hline
$10^{-1}$ & $0{,}497363752535$ & $7{,}947\times10^{-2}$ & $0{,}539402252170$ & $1{,}666\times10^{-3}$\\
$10^{-2}$ & $0{,}536085981012$ & $7{,}804\times10^{-3}$ & $0{,}540293300875$ & $1{,}667\times10^{-5}$\\
$10^{-4}$ & $0{,}540260231419$ & $7{,}787\times10^{-5}$ & $0{,}540302304968$ & $1{,}667\times10^{-9}$\\
$10^{-6}$ & $0{,}540301885121$ & $7{,}787\times10^{-7}$ & $0{,}540302305896$ & $5{,}130\times10^{-11}$\\
$10^{-8}$ & $0{,}540302302898$ & $5{,}497\times10^{-9}$ & $0{,}540302308449$ & $4{,}777\times10^{-9}$\\
$10^{-10}$ & $0{,}540302247387$ & $1{,}082\times10^{-7}$ & $0{,}540302247387$ & $1{,}082\times10^{-7}$\\
\hline
\end{tabular}
\end{center}

Pola galatnya bukan kebetulan. Untuk beda maju, ekspansi Taylor memberi
\[
  \frac{f(x_{0}+h)-f(x_{0})}{h}=f'(x_{0})+\frac{h}{2}f''(x_{0})+\mathcal{O}(h^{2}),
\]
jadi galat relatifnya berorde $h$ dan turun sepuluh kali setiap $h$ dibagi sepuluh:
$7{,}947\times10^{-2}$, $7{,}804\times10^{-3}$, $7{,}787\times10^{-5}$. Angkanya dapat
dicocokkan: $\frac{h}{2}f''(1)=-\frac{1}{2}\cdot10^{-2}\sin 1=-4{,}207\times10^{-3}$, dan
$4{,}207\times10^{-3}/0{,}5403=7{,}787\times10^{-3}$. Beda pusat memberi
\[
  \frac{f(x_{0}+h)-f(x_{0}-h)}{2h}=f'(x_{0})+\frac{h^{2}}{6}f'''(x_{0})+\mathcal{O}(h^{4}),
\]
sehingga galatnya berorde $h^{2}$ dan turun seratus kali: $1{,}666\times10^{-3}$,
$1{,}667\times10^{-5}$, $1{,}667\times10^{-9}$. Di sini $\frac{10^{-4}}{6}\cos 1
=9{,}005\times10^{-6}$, dan $9{,}005\times10^{-6}/0{,}5403$ kembali memberi $1{,}667\times10^{-5}$.

Pada $h=10^{-8}$ dan $10^{-10}$ beda pusat memburuk lagi: galat pemotongan sudah lebih kecil
daripada galat pembulatan yang tumbuh seperti $\varepsilon/h$, dengan
$\varepsilon\approx2{,}2\times10^{-16}$. Yang terbaik di tabel, $h=10^{-6}$ dengan galat
$5\times10^{-11}$, ordenya sesuai $\varepsilon^{1/3}\approx6\times10^{-6}$.

Kedua, gradien $f(x,y)=x^{2}y^{3}$ di $(2,-1)$ dengan nilai analitik $(-4,12)$. Beda pusat
memberi galat $3{,}333\times10^{-5}$ pada komponen $y$ untuk $h=10^{-2}$, turun ke
$3{,}333\times10^{-9}$ untuk $h=10^{-4}$: orde $h^{2}$. Komponen $x$ tepat sampai galat mesin,
sebab $f$ kuadratik dalam $x$.

Ketiga, aturan rantai pada komposisi berdimensi berbeda. Misalkan
\[
  H(x,y)=\bigl(x^{2}+y,\;xy,\;y^{3}\bigr),\qquad
  G(u,v,w)=\bigl(uv+w,\;u^{2}-vw\bigr),
\]
dan tinjau $F=G\circ H$ di $a=(1,2)$. Di sini $H(a)=(3,2,8)$, dan
\[
  J_{H}(1,2)=\begin{pmatrix}2&1\\ 2&1\\ 0&12\end{pmatrix}_{3\times2},
  \qquad
  J_{G}(3,2,8)=\begin{pmatrix}2&3&1\\ 6&-8&-2\end{pmatrix}_{2\times3}.
\]
Hasil kalinya berukuran $2\times2$:
\[
  J_{G}(3,2,8)\,J_{H}(1,2)=
  \begin{pmatrix}2&3&1\\ 6&-8&-2\end{pmatrix}
  \begin{pmatrix}2&1\\ 2&1\\ 0&12\end{pmatrix}
  =\begin{pmatrix}10&17\\ -4&-26\end{pmatrix}.
\]
\kertas{Untuk komposisi di atas, hitung $J_{F}$ di $(2,1)$ dengan dua jalan: hasil kali
$J_{G}(H(2,1))J_{H}(2,1)$ dan turunan langsung bentuk eksplisit. Lalu periksa arah $v=(1,3)$
dengan membandingkan $J_{F}(2,1)v$ terhadap $\bigl(F(a+sv)-F(a)\bigr)/s$ untuk $s=10^{-2}$
dan $s=10^{-4}$. Mengapa galatnya turun hanya sebesar faktor $10$, bukan $100$?}

\section{Lab komputasi}
\tingkat{lanjut}

Kode berikut dijalankan apa adanya dengan numpy 2.5.2 pada Python 3: bagian pertama mengulang
tiga hitungan di atas, lalu satu percobaan yang tidak dapat dikerjakan dengan tangan, beda
hingga di bawah derau pengukuran.

\begin{lab}
\begin{verbatim}
import numpy as np

# 1. beda hingga untuk f(x) = sin x di x0 = 1
x0 = 1.0
eksak = float(np.cos(x0))
print("cos(1) =", eksak)
for h in [1e-1, 1e-2, 1e-4, 1e-6, 1e-8, 1e-10]:
    maju = (np.sin(x0 + h) - np.sin(x0)) / h
    pusat = (np.sin(x0 + h) - np.sin(x0 - h)) / (2*h)
    em = abs(maju - eksak) / abs(eksak)
    ep = abs(pusat - eksak) / abs(eksak)
    print(f"{h:8.0e} fwd={maju:.12f} {em:.3e} "
          f"cen={pusat:.12f} {ep:.3e}")
\end{verbatim}
Keluaran sebenarnya:
\begin{verbatim}
cos(1) = 0.5403023058681398
   1e-01 fwd=0.497363752535 7.947e-02 cen=0.539402252170 1.666e-03
   1e-02 fwd=0.536085981012 7.804e-03 cen=0.540293300875 1.667e-05
   1e-04 fwd=0.540260231419 7.787e-05 cen=0.540302304968 1.667e-09
   1e-06 fwd=0.540301885121 7.787e-07 cen=0.540302305896 5.130e-11
   1e-08 fwd=0.540302302898 5.497e-09 cen=0.540302308449 4.777e-09
   1e-10 fwd=0.540302247387 1.082e-07 cen=0.540302247387 1.082e-07
\end{verbatim}
\end{lab}

Baris terakhir layak dibaca dua kali: pada $h=10^{-10}$ kedua rumus memberi angka identik,
karena keduanya menghitung galat pembulatan yang sama, bukan turunan.

\begin{lab}
\begin{verbatim}
def f(p):
    return p[0]**2 * p[1]**3

def grad_fd(fun, a, h):
    g = np.zeros_like(a)
    for j in range(a.size):
        e = np.zeros_like(a)
        e[j] = h
        g[j] = (fun(a + e) - fun(a - e)) / (2*h)
    return g

a = np.array([2.0, -1.0])
g_eks = np.array([2*a[0]*a[1]**3, 3*a[0]**2*a[1]**2])
print("gradien =", g_eks)
for h in [1e-2, 1e-4, 1e-6]:
    g = grad_fd(f, a, h)
    print(f"{h:8.0e} {g[0]:.12f} {g[1]:.12f} "
          f"{abs(g[0]+4)/4:.3e} {abs(g[1]-12)/12:.3e}")
\end{verbatim}
Keluaran sebenarnya:
\begin{verbatim}
gradien = [-4. 12.]
   1e-02 -4.000000000000 12.000400000000 1.577e-14 3.333e-05
   1e-04 -4.000000000004 12.000000039998 1.000e-12 3.333e-09
   1e-06 -4.000000000115 11.999999999679 2.876e-11 2.676e-11
\end{verbatim}
\end{lab}

\begin{lab}
\begin{verbatim}
def hh(p):
    x, y = p
    return np.array([x*x + y, x*y, y**3])

def gg(q):
    u, v, w = q
    return np.array([u*v + w, u*u - v*w])

def F(p):
    return gg(hh(p))

def jac(fun, p, step):
    n = p.size
    J = np.zeros((fun(p).size, n))
    for j in range(n):
        e = np.zeros(n)
        e[j] = step
        J[:, j] = (fun(p + e) - fun(p - e)) / (2*step)
    return J

a = np.array([1.0, 2.0])
q = hh(a)
Jh = np.array([[2*a[0], 1.0], [a[1], a[0]], [0.0, 3*a[1]**2]])
Jg = np.array([[q[1], q[0], 1.0], [2*q[0], -q[2], -q[1]]])
J_rantai = Jg @ Jh
J_num = jac(F, a, 1e-6)
print("h(1,2) =", q)
print("Jg @ Jh =\n", J_rantai)
print("selisih maks =", f"{np.max(np.abs(J_num - J_rantai)):.3e}")

v = np.array([0.7, -1.3])
d_rantai = J_rantai @ v
d_num = (F(a + 1e-6*v) - F(a - 1e-6*v)) / 2e-6
print("J v =", d_rantai)
print("beda hingga =", d_num)
print("galat relatif =",
      f"{np.max(np.abs(d_num - d_rantai)/np.abs(d_rantai)):.3e}")
\end{verbatim}
Keluaran sebenarnya:
\begin{verbatim}
h(1,2) = [3. 2. 8.]
Jg @ Jh =
 [[ 10.  17.]
 [ -4. -26.]]
selisih maks = 1.858e-09
J v = [-15.1  31. ]
beda hingga = [-15.1  31. ]
galat relatif = 9.860e-11
\end{verbatim}
\end{lab}

Selisih $1{,}858\times10^{-9}$ antara Jacobian aturan rantai dan Jacobian beda hingga adalah
galat pemotongan beda pusat, bukan kesalahan aljabar; pemeriksaan arah memberi
$9{,}86\times10^{-11}$, sebab keduanya memakai titik yang sama.

\begin{lab}
\begin{verbatim}
rng = np.random.default_rng(7)
sigma = 1e-3
for h in [1e-1, 1e-2, 1e-3, 1e-4, 1e-5]:
    galat = []
    for _ in range(2000):
        s1 = np.sin(x0 + h) + sigma*rng.standard_normal()
        s0 = np.sin(x0) + sigma*rng.standard_normal()
        galat.append(abs((s1 - s0)/h - eksak))
    print(f"{h:8.0e} galat rerata = {np.mean(galat):.3e}")
\end{verbatim}
Keluaran sebenarnya:
\begin{verbatim}
   1e-01 galat rerata = 4.203e-02
   1e-02 galat rerata = 1.112e-01
   1e-03 galat rerata = 1.160e+00
   1e-04 galat rerata = 1.117e+01
   1e-05 galat rerata = 1.088e+02
\end{verbatim}
\end{lab}

Percobaan terakhir adalah inti bab ini dalam angka. Dengan derau
$\sigma=10^{-3}$, hasil bagi membawa dua galat: bias pemotongan $\approx0{,}42h$ dan derau
yang menguat seperti $\sqrt{2}\sigma/h$. Pada $h=10^{-1}$ bias mendominasi ($0{,}042$), pada
$h=10^{-4}$ derau mendominasi ($11{,}2$), dan keduanya seimbang ketika
$0{,}42h=\sqrt{2}\cdot10^{-3}/h$, yaitu $h\approx0{,}058$: dengan alat seburuk itu, langkah
terbaik sekitar enam perseratus, bukan seperseratus juta.

Inilah jawaban jujur untuk pertanyaan pembuka. Turunan adalah limit, tetapi aproksimasi
terbaiknya bukan limit yang sedekat mungkin. Yang terbaik adalah kompromi antara dua galat
yang bergerak berlawanan arah, dan kompromi itu dapat dihitung sebelum satu angka pun
diukur.

\begin{untukdosen}
Urutan yang berhasil dipakai: definisi hasil bagi, lalu bentuk aproksimasi linear, baru
aturan hitung. Bentuk itu terasa berlebihan untuk satu variabel, tetapi hanya ia yang bertahan
sampai Jacobian; bila diperkenalkan belakangan, mahasiswa menghafal aturan rantai sebagai
perkalian matriks tanpa tahu mengapa urutannya begitu.

Dua miskonsepsi khas. Pertama, ``turunan parsial ada berarti fungsi mulus''; contoh
$xy/\sqrt{x^{2}+y^{2}}$ di papan mematahkannya asalkan dihitung bersama: parsial nol,
kontinu, tetapi arah diagonal memberi $1/2$. Kedua, arah perkalian pada aturan rantai:
mahasiswa menulis $J_{H}J_{G}$ karena urutan penulisan fungsi terbalik dari urutan
pengerjaan; cara memeriksanya adalah ukuran matriksnya.

Aktivitas papan tulis: minta kelas menghitung kedua parsial di titik asal untuk kedua contoh
penyangkal, lalu memutuskan apakah fungsi itu punya bidang singgung. Tanda kesiapan lanjut:
mahasiswa menjelaskan mengapa $D_{u}f=\nabla f\cdot u$ gagal pada contoh kedua, dan menyebut
kata ``sisa'', bukan ``rumus''.
\end{untukdosen}

\section{Ringkasan}
\tingkat{dasar}

\begin{ringkasan}
\begin{itemize}
  \item Turunan $f'(c)$ adalah limit hasil bagi selisih, sekaligus koefisien tunggal
        aproksimasi linear $f(c+h)=f(c)+Lh+r(h)$ dengan $r(h)/h\to0$. Diferensiabilitas
        menjamin kekontinuan, bukan sebaliknya.
  \item Aturan rantai $(f\circ g)'=f'(g(c))g'(c)$ dibuktikan lewat aproksimasi linear, dan
        buktinya bertahan ke dimensi banyak: $J_{G\circ H}(a)=J_{G}(H(a))J_{H}(a)$.
  \item Teorema nilai rata-rata menghubungkan tanda turunan dengan perilaku fungsi: $f'\ge0$
        berarti tidak turun, $|f'|\le M$ berarti Lipschitz; ia gagal untuk fungsi vektor.
  \item Taylor dengan sisa Lagrange memakai $\xi$ yang tidak diketahui, sehingga yang dipakai
        hanya batas atas $|f^{(n)}|$.
  \item Untuk banyak variabel, turunan adalah peta linear dengan matriks Jacobian. Parsial
        boleh ada tanpa kekontinuan, dan boleh ada bersama kekontinuan tanpa
        diferensiabilitas.
  \item Gradien menunjuk arah naik tercepat dengan laju $\|\nabla f\|$, tegak lurus kontur;
        Hessian yang definit menentukan jenis titik kritis.
  \item Beda hingga berorde $h$ (maju) atau $h^{2}$ (pusat), dan galat pembulatan tumbuh
        seperti $\varepsilon/h$. Langkah terkecil bukan langkah terbaik.
\end{itemize}
\end{ringkasan}

\section{Latihan}
\tingkat{menengah}

\begin{latihan}
\begin{enumerate}
  \item Hitung $f'(2)$ dari definisi untuk $f(x)=x^{3}$ dan $f(x)=\sqrt{x}$, sampai faktor
        $h$ hilang dari penyebut.

  \item Buktikan bahwa bila $f'(c)>0$, ada $\delta>0$ sehingga $f(c+h)>f(c)$ untuk setiap
        $h\in(0,\delta)$, dan jelaskan mengapa ini tidak berarti $f$ naik pada suatu interval.

  \item Buktikan aturan hasil bagi dengan menurunkan $1/g$ lebih dulu, lalu memakai aturan
        hasil kali. Sebutkan di langkah mana kekontinuan $g$ dipakai.

  \item Nilai klaim ini: \emph{bila $f$ mempunyai semua turunan parsial di suatu titik, maka
        $f$ kontinu di titik itu.} Nyatakan benar, salah, atau belum cukup bukti, dan
        sertakan contoh penyangkal bila perlu.

  \item Untuk $f(x,y)=x^{3}y-xy^{3}$, hitung gradien di $(2,1)$, tentukan arah naik tercepat,
        dan hitung laju maksimumnya. Periksa dengan beda pusat $h=10^{-4}$ dan bandingkan
        galat relatifnya dengan orde $h^{2}$.

  \item Buktikan bahwa bila $f'=0$ di mana-mana pada suatu interval, maka $f$ konstan.
        Jelaskan mengapa bukti itu gagal pada himpunan yang tidak terhubung, dan beri
        contohnya.

  \item Misalkan $H(x,y)=(e^{x}y,\;x+y^{2})$ dan $G(u,v)=(u^{2}v,\;u-v)$. Hitung
        $J_{G\circ H}$ di $(0,1)$ dengan dua jalan: hasil kali matriks dan turunan langsung.
        (Petunjuk: $H(0,1)=(1,1)$.)

  \item Tentukan semua titik kritis $f(x,y)=x^{3}+y^{3}-3xy$, lalu klasifikasikan dengan
        Hessian; jelaskan mengapa tanda determinan tidak cukup bila determinannya nol.

  \item Taksir $\sqrt{1{,}04}$ dengan polinomial Taylor orde dua di $a=1$, lalu batas
        galatnya dengan sisa Lagrange. Bandingkan batas itu dengan galat sebenarnya: apakah batasnya ketat? (Petunjuk:
        $f''(x)=-\frac14x^{-3/2}$ dan $f'''$ positif pada selang itu.)
\end{enumerate}
\end{latihan}

\begin{untukpeneliti}
Yang terbukti di bab ini adalah pernyataan klasik satu dan banyak variabel: aproksimasi
linear dan ketunggalan petanya, aturan rantai, teorema nilai rata-rata dan akibatnya, Taylor
dengan sisa Lagrange, kesetaraan diferensiabilitas dengan Jacobian, dan syarat cukup
kekontinuan parsial. Semua itu teorema, bukan konvensi, berlaku di ruang berdimensi hingga
atas $\mathbb{R}$.

Yang sering dipopulerkan adalah klaim ``gradien adalah turunan''. Tanpa diferensiabilitas
klaim itu salah, dan $xy/\sqrt{x^{2}+y^{2}}$ adalah penyangkal yang tetap kontinu. Klaim
``Hessian menentukan jenis titik kritis'' hanya berlaku bila Hessian nonsingular; kasus
singular memerlukan turunan orde lebih tinggi atau analisis langsung.

Sumber primer. Rudin, \emph{Principles of Mathematical Analysis}, bab
\emph{Differentiation} dan \emph{Functions of Several Variables}, memuat seluruh pernyataan
di atas. Spivak, \emph{Calculus on Manifolds}, bab \emph{Differentiation}, mengambil jalan
yang sama dan lebih singkat. Stewart, \emph{Calculus: Early Transcendentals} edisi ke-9,
bagian 14.3 sampai 14.7, memberi padanan komputasional, termasuk teorema Clairaut.

Batas keberlakuan: semua argumen menganggap himpunan definisi terbuka dan turunan yang
dipakai ada. Pada fungsi yang hanya terdiferensialkan hampir di mana-mana, kesimpulan
monotonisitas dan Taylor runtuh, dan itu wilayah teori ukur yang tidak dibahas buku ini.
Untuk fungsi vektor, teorema nilai rata-rata tidak berlaku; penggantinya ketaksamaan nilai
rata-rata yang menuntut norma operator. Orde beda hingga terbaik bergantung pada skala
fungsi, derau, dan aritmetika; bab ini hanya menunjukkan ordenya secara empiris.
\end{untukpeneliti}

\peta{Bab ini mengikuti Rudin, \emph{Principles of Mathematical Analysis}, bab
\emph{Differentiation} dan \emph{Functions of Several Variables}; versi komputasionalnya
Stewart, \emph{Calculus: Early Transcendentals} edisi ke-9, bagian 14.3 sampai 14.7; padanan
abstraknya Spivak, \emph{Calculus on Manifolds}, bab \emph{Differentiation}. Kekontinuan dan
teorema nilai ekstrem dipakai dari Bab 11, ketaksamaan Cauchy--Schwarz dari Bab 10, norma
Jacobian kembali di Bab 23 dan Bab 28.}
